% All-degree characteristic-two contact argument; independent audit requested.
\subsection{Contact curves in every degree in characteristic two}
\label{subsec:all-degree-char2-contact}

The following argument removes the degree restriction from the
contact-incidence method.  Its global dimension input is
Christ--He--Tyomkin (2023), Proposition~2.7 (p.~8) and
Corollary~5.3 (p.~27), in the corrected publisher version of
\emph{On the Severi problem in arbitrary characteristic}.
Both statements apply over every algebraically closed field,
including characteristic two.  If a reduced plane curve of degree
$s$ has $c$ normalization components of genera $h_1,\ldots,h_c$,
their geometric genus convention is $1-c+\sum h_i$.  Consequently
the locus of such curves has dimension at most
\begin{equation}\label{eq:contact-severi-bound}
 3s+\sum_i h_i-c.
\end{equation}
If a family attains this dimension, its generic member is nodal.
The publisher's correction concerns the article history and missing
references; the opened corrected text contains these statements.

\begin{lemma}[The conductor inclusion for an even contact]
\label{lem:contact-conductor}
Let $C$ be a smooth plane curve of degree $d$ over an algebraically
closed field of characteristic two, and let $Q$ be a reduced plane
curve of degree $s<d$ with $Q\cap C=2D$ as divisors on $C$.
Write $\nu:\widetilde Q=\coprod Q_i\to Q$ for the normalization,
and write
\[
 \operatorname{div}(f|_{Q_i})=2A_i+\epsilon_i,
 \qquad M_i=\mathcal O_{Q_i}(A_i),
\]
where $f$ defines $C$ and $\epsilon_i$ is the reduced odd-order
divisor.  It is supported above the singular points of $Q$.
Set
\[
 J=\mathcal I_{D/Q}(s),\qquad
 L_C=\mathcal O_C(D),\qquad
 b=d-s.
\]
Then
\begin{equation}\label{eq:general-contact-reciprocity}
 h^0(C,L_C)=1+h^0(Q,J),
\end{equation}
and there is an inclusion
\begin{equation}\label{eq:conductor-section-inclusion}
 \nu_*\left(\bigoplus_i
       \omega_{Q_i}\otimes M_i^{-1}(3)\right)\hookrightarrow J.
\end{equation}
In particular, with $t=\sum_i h^1(Q_i,M_i)$,
\begin{equation}\label{eq:contact-h1-cancellation}
 h^0(Q,J)\ge t,
 \qquad t>0\Longrightarrow h^0(Q,J)\ge t+1.
\end{equation}
\end{lemma}

\begin{proof}
The same subscheme $D\subset\mathbf P^2$ is used on both curves.
Its ideal-sheaf sequences, twisted by $s$, are
\[
 0\to\mathcal O_{\mathbf P^2}(s-d)
  \to\mathcal I_D(s)\to L_C\to0,
 \qquad
 0\to\mathcal O_{\mathbf P^2}
  \to\mathcal I_D(s)\to J\to0.
\]
Here $\mathcal O_C(s)(-D)\simeq\mathcal O_C(D)$ because
$2D=Q|C$.  Since $s<d$ and the first cohomology of every line
bundle on $\mathbf P^2$ vanishes, these sequences prove
\eqref{eq:general-contact-reciprocity}.

We prove the inclusion locally and then express it globally.
At a point of $D$, complete the local ring of the plane.  Since
$C$ is smooth, choose its equation $f$ as one parameter and let
$b_0$ cut out the half divisor on $C$.  After multiplying the local
equation of $Q$ by a unit, that equation is
$b_0^2-fa_0$.  In the reduced one-dimensional ring
$A=\widehat{\mathcal O}_{Q,P}$, the ideal of $D$ is
$I=(f,b_0)$ and $f$ is a non-zero-divisor.
Reduction modulo $f$ gives
$A/(f)\simeq k[[x]]/(b_0^2)$, where the annihilator of $b_0$ is
$(b_0)$.  Hence
\[
 I^*:=\operatorname{Hom}_A(I,A)
 =\{z:zI\subset A\}=f^{-1}I,
 \qquad I^{**}=I.
\]
These equalities can be read directly in the total quotient ring:
$zf\in A$ and $zb_0\in A$ say that $zf$ belongs to
$(f):(b_0)=I$.

Let $v_i$ be the valuation on a normalization branch.  Put
\[
 r_i=v_i(f),\qquad j_i=\min\{v_i(f),v_i(b_0)\},
 \qquad\gamma_i=2j_i-r_i.
\]
The relation $b_0^2=fa_0$ gives $\gamma_i\ge0$; moreover
$\gamma_i\equiv r_i\pmod2$.  Thus
$\gamma_i=\epsilon_i(P_i)+2k_i$ for an integer $k_i\ge0$.
Write $c_i$ for the conductor exponent on this branch, so that
the conductor ideal on the normalization is generated there by
a parameter to the power $c_i$.
If a tuple $z$ of branch functions satisfies
\[
 v_i(z)\ge c_i+r_i-j_i\quad\hbox{for every branch},
\]
then its product with every member of $I^*=I/f$ lies in the
conductor, hence in $A$.  It follows that $z\in I^{**}=I$.
No separate assumption on a tame singularity, on its number of
branches, or on its analytic type is used here.

Let $\Delta$ be the conductor divisor on $\widetilde Q$, let
$\Gamma=\sum\gamma_iP_i$, and let $K_0=\sum k_iP_i$.
If $N=\nu^*J/\text{torsion}$, the valuation calculation gives
\[
 N|_{Q_i}\simeq M_i(-b)(-K_0|_{Q_i}),\qquad
 \nu_*N(\Gamma-\Delta)\subset J.
\]
Indeed the normalized ideal has vanishing order $j_i$, and
$\gamma_i=2j_i-r_i$ is precisely the adjustment from its square
to the divisor of $f$.

For a reduced plane curve, finite duality (Stacks Project,
Remark~48.20.12, Tag~0AX4, opened and checked) and adjunction give
\[
 \omega_{\widetilde Q}
 \simeq\nu^*\omega_Q(-\Delta)
 =\mathcal O_{\widetilde Q}(s-3)(-\Delta).
\]
For completeness, the conductor occurs because
$\operatorname{Hom}_{\mathcal O_Q}
(\nu_*\mathcal O_{\widetilde Q},\mathcal O_Q)$ is the conductor
ideal; finite duality identifies its tensor product with
$\omega_Q$ with $\nu_*\omega_{\widetilde Q}$.
Using $M_i^2\simeq\mathcal O_{Q_i}(d)(-\epsilon_i)$ now yields
\[
 \omega_{Q_i}\otimes M_i^{-1}(3)
 \simeq M_i(-b)(\epsilon_i-\Delta_i)
 =N|_{Q_i}(K_{0,i}+\epsilon_i-\Delta_i).
\]
Since $\Gamma=\epsilon+2K_0$, this line bundle is a subsheaf
of $N(\Gamma-\Delta)$, proving
\eqref{eq:conductor-section-inclusion}.

By Serre duality, $h^1(M_i)=h^0(\omega_{Q_i}M_i^{-1})$.
Twisting by the positive line bundle $\mathcal O_{Q_i}(3)$
does not decrease this dimension.  If the initial space is nonzero,
the increase is strict: choose a point where some section does not
vanish, and choose a plane cubic through its image whose pullback
is a nonzero divisor.  Multiplication by that cubic has image
consisting of sections vanishing there, whereas multiplying the
chosen initial section by a cubic not through the point gives a
section outside the image.  Summing over the components proves
\eqref{eq:contact-h1-cancellation}.
\end{proof}

\begin{lemma}[No lower-degree reduced contact for a generic curve]
\label{lem:no-low-degree-contact}
Fix $d\ge4$ and $0<s<d-3$ in characteristic two.
A generic smooth plane curve $C$ of degree $d$ has no reduced
plane curve $Q$ of degree $s$ with $Q|C$ even.
\end{lemma}

\begin{proof}
Let $N_d=\binom{d+2}{2}-1$.
For fixed normalization genera, component degrees, and odd-parity
divisors $\epsilon_i$, the line bundles
$M_i^2=\mathcal O_{Q_i}(d)(-\epsilon_i)$ have finite choices.
They can be parameterized after a finite base change of the
singularity stratum; inseparable base changes do not affect
dimension.  The normalized restrictions of $f$ are squares of
sections of these $M_i$, times the fixed odd-parity factors.
Ignoring descent conditions only increases the dimension.
Riemann--Roch gives
\[
 \sum_i h^0(M_i)
 =\frac{ds}{2}-\frac{|\epsilon|}{2}
   +c-\sum_i h_i+t,
 \qquad t=\sum_i h^1(M_i).
\]
The kernel of restricting degree-$d$ forms to $Q$ has dimension
$\binom{d-s+2}{2}$.  Combining this with
\eqref{eq:contact-severi-bound}, and projectivizing, gives the
incidence bound
\begin{equation}\label{eq:general-contact-dimension}
 \dim\mathcal I_{d,s}
 \le N_d-\frac{s(d-s-3)}2-\frac{|\epsilon|}2+t
\end{equation}
on each stratum with fixed $t$.

Here is the precise fibre argument used with this bound.
Suppose a component of the even-contact incidence dominates the
ordinary smooth degree-$d$ parameter space.  Over its generic
curve, roots $L_C^2=\mathcal O_C(s)$ form a finite scheme.
For each geometric root the contact forms are the projective
space of its nonzero sections, followed by squaring.
Restriction of degree-$s$ plane forms to $C$ is an isomorphism,
so no further choices of $Q$ occur.  The reduced-contact condition
is open in this projective section space.  The generic fibre of
our component therefore has dimension
$h^0(L_C)-1=h^0(J)$.
Its generic point belongs to a single stratum of the finite
stratification by the data above, and an open subset of the same
generic fibre remains in that stratum.  Thus the fibre dimension
may be subtracted in \eqref{eq:general-contact-dimension}; we are
not subtracting a family of contact curves from a special stratum
that it need not preserve.

The conductor lemma gives $h^0(J)\ge t$.  Dominance would imply
\[
 N_d\le N_d-\frac{s(d-s-3)}2-\frac{|\epsilon|}2<N_d,
\]
a contradiction.  There are finitely many relevant component
degrees, genera, parity choices in finite covers, and cohomology
dimensions.  Thus their images have proper closure, proving the
claim for a generic curve.
\end{proof}

\begin{theorem}[The theta orbits of generic plane curves in characteristic two]
\label{thm:all-char2-theta-orbits}
Let $d\ge4$, $g=(d-1)(d-2)/2$, and let $C$ be the geometric
generic ordinary smooth plane curve of degree $d$ in characteristic
two.
If $d$ is even, every theta characteristic other than the canonical
one $B_C$ has exactly one section, and geometric monodromy is
transitive on these $2^g-1$ characteristics.
If $d$ is odd, every theta characteristic other than $B_C$ and
$\mathcal O_C((d-3)/2)$ has exactly one section, and geometric
monodromy is transitive on these $2^g-2$ characteristics.
Equivalently, on $V^\vee$ geometric monodromy is transitive on
all nonzero vectors in even degree, and on
$V^\vee\setminus\{0,v\}$ in odd degree, where $v$ is the
distinguished nonzero invariant vector.
\end{theorem}

\begin{proof}
Put $e=d-3$.  A section of an effective theta characteristic
$\theta_C$ gives a contact form $Q$ of degree $e$, and conversely
such a form gives an effective theta characteristic and its
section up to scalar.  Work throughout over the ordinary smooth
degree-$d$ locus, so theta characteristics have finitely many
geometric choices, with constant count $2^g$.  Passing to a
purely inseparable cover if necessary changes neither dimensions
nor geometric Galois orbits.

\paragraph{The unique dominant smooth-contact incidence.}
For a smooth $Q$, contact reciprocity gives
$h^0(\theta_C)=1+h^0(\theta_Q)$, where
$\theta_Q=\mathcal O_Q(D-3)$.
Riemann--Roch gives
\[
 h^0(\theta_Q(3))=3e+h^0(\theta_Q(-3)).
\]
Thus the incidence dimension, for a stratum with these dimensions
constant, is at most $N_d+h^0(\theta_Q(-3))$.
If $\theta_Q$ is effective, then
$h^0(\theta_Q(-3))<h^0(\theta_Q)$: multiplication by a suitable
cubic gives a proper subspace, as in the final paragraph of the
conductor lemma.  The fibre dimension $h^0(\theta_C)-1$
is $h^0(\theta_Q)$, so such an incidence cannot dominate.
For an ordinary smooth $Q$, its canonical $B_Q$ is the unique
non-effective theta characteristic.  It has
$h^0(B_Q(3))=3e$.  For a nonordinary smooth $Q$, all its geometric
theta characteristics are effective.  Consequently the only
possible dominant smooth-contact incidence is the canonical one
over ordinary smooth $Q$.

That canonical incidence is geometrically irreducible of dimension
$N_d$: the ordinary degree-$e$ base is geometrically irreducible;
after Frobenius base change the relative canonical theta
characteristic exists; its $3e$ sections form a vector bundle;
and lifts of their squares through restriction form affine spaces
with the constant ten-dimensional kernel
$Q\cdot H^0(\mathbf P^2,\mathcal O(3))$.
Projectivization gives total dimension
$\binom{e+2}{2}-1+3e+10-1=N_d$.
For $e=1,2$ the genus-zero canonical theta characteristic is used,
and the same computation applies.

\paragraph{Reduced singular contact forms cannot dominate.}
Apply the notation and incidence bound of the preceding lemmas
with $s=e$.  A hypothetical dominating component has dimension
at most $N_d-|\epsilon|/2+t$, and generic fibre dimension
$h^0(J)$.  If $t>0$, the strict part of
\eqref{eq:contact-h1-cancellation} gives a contradiction.
If $|\epsilon|>0$, its total degree is even and at least two,
so the non-strict inequality also gives a contradiction.
If $h^0(J)>0$ while $t=|\epsilon|=0$, the fibre dimension again
gives a contradiction.  Only
\[
 t=0,\qquad \epsilon=0,\qquad h^0(J)=0
\]
could remain for a dominant component.

In this case its image in the degree-$e$ curve parameter space
must attain the maximal Severi dimension
$3e+\sum h_i-c$; otherwise the incidence bound loses at least
one further dimension.  By Christ--He--Tyomkin Corollary~5.3,
the generic $Q$ is therefore nodal.  Since it is singular, choose
one node.  The conductor inclusion and $h^0(J)=0$ imply
\[
 H^0(Q_i,\omega_{Q_i}M_i^{-1}(3))=0
 \quad\hbox{on every component}.
\]
If the two preimages $P_+,P_-$ of the node lie on one component,
there exists a plane cubic whose pullback contains both points.
It follows that
$H^0(\omega M_i^{-1}(P_++P_-))=0$, hence
$H^1(M_i(-P_+-P_-))=0$.  The evaluation map at the two points
is therefore onto.  If they lie on different components, the
same argument with one point proves surjectivity of each
one-point evaluation.  In either case, equality of the two values
needed for descent of the squared restriction imposes one
nonzero linear condition on the root sections after taking square
roots of its scalar coefficients.  This lowers the incidence
dimension by at least one, the final contradiction.

\paragraph{Nonreduced contact forms.}
Every nonreduced form of degree $e$ has a unique expression up to
scalars $Q=R^2S$, with $S$ reduced and $\deg S=s<e$.
If $s>0$, evenness of $Q|C$ implies evenness of $S|C$, excluded
for a generic $C$ by Lemma~\ref{lem:no-low-degree-contact}.
If $s=0$, then $e$ is even, and the half divisor of $Q|C$ is
the divisor of $R|C$.  It gives exactly the plane theta
characteristic $\mathcal O_C(e/2)$.  This case occurs precisely
when $d$ is odd and is the sole nonreduced contact case retained.

\paragraph{Counting the generic fibre and obtaining the orbit.}
Every noncanonical theta characteristic is effective.
All except the plane characteristic in odd degree must therefore
occur in the geometrically irreducible canonical smooth-contact
incidence.  This proves its dominance.  Since it has dimension
$N_d$, its generic fibre is finite, so each of these theta
characteristics has exactly one section up to scalar.
The generic geometric fibre has respectively $2^g-1$ or
$2^g-2$ points.  Irreducibility makes geometric Galois monodromy
transitive on them; purely inseparable factors of the incidence
do not change this orbit statement.  Translating theta
characteristics by the canonical $B_C$ identifies their set
with $V^\vee$, proving the asserted vector orbits.
\end{proof}
